% \chapter*{Executive Summary}
% \addcontentsline{toc}{chapter}{Executive Summary}
% \markboth{Executive Summary}{Executive Summary}
% \noindent Summarize the problem, the proposed solution, major design decisions, prototype status, evaluation results, and expected impact. Keep this section readable by sponsors and reviewers who may not read the full report.

\part{Project Definition and Plan}

\chapter{Team Declaration}
\ChapterLabel{TeamDeclaration}

\section{Team Descriptions}
Team Dynamo is a four-person senior design team, each majoring in Software Engineering. Throughout their academic careers, they have developed experience in areas such as software development, system design, testing, and problem-solving. As a team, they bring together individual skills, perspectives, and technical strengths to collaborate on the design and development of a complete software solution. Their senior design project gives them the opportunity to apply what we have learned in the classroom to a practical challenge while strengthening our abilities in teamwork, communication, project management, and the software development process. The following paragraphs introduce each team member and highlight their individual experience and interests.

%Jacky
\subsection{Jacky Lei}
A senior at Stevens Institute of Technology pursuing a Bachelor’s of Engineering in Software Engineering and Mathematics and Physics minors. After graduation, I hope to be enrolled in a PhD program researching computational astrophysical magnetohydrodynamics. I am a part of research labs in optics and computational fluid dynamics. Outside of academics I love to run, perform wushu, and cook.

%Kayla
\subsection{Kayla Holmes}
Kayla Holmes is a Senior Software Engineer studying at Stevens Institute of technology. I am passionate about systems and optimization of code. In my free time I enjoy doing arts and crafts and reading. I am a part of various culture organizations as well as the theater club as the Vice President.

%Robert
\subsection{Robert Galletta}
A Stevens Institute of Technology senior with a Software Engineering major and a minor in Ocean Engineering. He works at the iSMART lab in the ABS center and maintains the landscaping of the local Chi Phi Fraternity chapter. 

%Nica
\subsection{Nica Saoi}
Nica Saoi is an undergraduate student at Stevens Institute of Technology, pursuing a major in software engineering, a minor in computer science, and a graduate certificate in software design and development. What really drew her to software engineering, and technology in general is the human element. She is interested in how people interact with technology, especially as artificial intelligence becomes an important aspect of our daily lives. She wants to learn more about how technology shapes behavior, influences decisions, and affects what people use every day.

\chapter{Project Plan}
\ChapterLabel{ProjectPlan}
\subsection{Mission Statement}
\ChapterLabel{MissionStatement}
We aim to create a framework that improves \gls{approachability}, traceability and \gls{reproducibility} in research.

\subsection{Key Drivers}\ChapterLabel{KeyDrivers}

Currently, the only research artifact easily accessible to the public is the research paper. This poses a challenge to \gls{reproducibility} since papers often have formatting restrictions imposed by academic journals that lead to only the core details of the research being present in the paper, leaving out smaller details that might be crucial being able to replicate results. \cite{christensen2018transparency} \cite{johnson2021evaluating} By improving \gls{reproducibility} in research, we aim to make it easier for other researchers to build on existing work as well as verify research results.

Creating a platform that hosts additional artifacts such as data, code, additional figures, videos, and experiment logs, allows researchers to more easily maintain robust documentation and clarity of their work. \cite{miyakawa_no_2020} \cite{randall2018irreproducibility} Having a separate research artifact in addition to the research paper gives researchers more space to elaborate on their work without concern for word count, formatting, or other limits imposed by academic journals that might prevent expository content from being present in the paper. \cite{christensen2018transparency} \cite{johnson2021evaluating} In addition, other artifacts allow more ways for readers to understand material. For example, there are many cases where complex processes described in papers can be replaced by brief visualizations. Experiment logs may also give more context to the experiments that were presented as well as present failed research directions. 

Although there are existing research information management systems, they tend to be inward facing, limited in scope, and are not widespread. In contrast, the framework shall be made with third-party researchers and the general public in mind. 


\subsection{Key Constraints}\ChapterLabel{KeyConstraints}

In order to be useful to the general research community, the framework shall be usable without a coding background. Although coding is an important skill in many fields, requiring coding may serve as a barrier for undergraduate researchers or fields where coding is not a major part of daily research tasks. We want to instill \gls{reproducibility} and documentation values early on and throughout the scientific community, and having the framework be usable without code is important for approachability. 

To maximize the spread of better \gls{reproducibility} and documentation practices, the framework shall be applicable to any scientific field. These values are important to the scientific community as a whole and the framework therefore needs to be adaptable for any field.

\chapter{Development Plan}
\ChapterLabel{DevelopmentPlan}

The Dynamo website will be developed using an iterative approach. Instead of building the entire system at once, the team will work on smaller features that can be designed, tested, and improved over time. This gives the team room to respond to stakeholder feedback and update requirements as the project develops.

\section{Development Approach}

Development of the first prototype will first focus on the main workflows that support Dynamo's goals of
improving the approachability, traceability, and reproducibility of research. This includes creating and discovering research projects, managing research artifacts, and connecting information such as publications, figures, datasets, and experiments.

To enable a smooth transition from the prototype to the final product, we will use a NoSQL database such as ScyllaDB. ScyllaDB will allow the project to be scalable while being flexible in the beginning stages of development due to the lack of a fixed schema. 

For the first prototype, we will have everything hosted locally in Docker containers. In following prototypes we will iteratively move to a more scalable configuration. 

User stories, requirements, and use cases will continue to be refined as the team receives more feedback. These development artifacts will be tracked such that implemented features can be linked back to the needs that
motivated them.


\section{Development Process}

As we follow an iterative approach, we will be rapidly moving through the V-model of software development. In each iteration, we will analyze, refine, and/or add requirements, update the architecture, modify/add code, test, and deploy. To automate movement through the V-model, we will incorporate some DevOps tools such as GitHub Actions, Git, and Docker. When possible, we aim to use a test-driven approach for development to ensure expectations of the code itself are clear. 

Team members will divide work across areas such as frontend development, backend development, database design, testing, documentation, and stakeholder research.

Progress will be coordinated through regular team meetings and tracked development tasks.

\section{Testing and Project Changes}

Testing will take place throughout development. Functional testing will be used to check that features satisfy their requirements and use cases, while integration testing will confirm that major system components work together correctly.

Changes to requirements or project scope will be reviewed based on their importance, technical impact, dependencies, and available development time. 

%\section{How to Use This Chapter}
%Describe governance, roles, milestones, change control, and risks. Adapt the plan to the lifecycle the team will actually follow.
%\section{Example}
%The product owner prioritizes outcomes; the architect owns significant design decisions; developers implement approved work; and the test lead retains evidence. Typical gates are needs reviewed, requirements baselined, architecture approved, evidence complete, and release approved. Example risk: an external service may be unavailable; define the failure response and test it before release.

%\chapter{Background}\ChapterLabel{Background}
%Explain the context, prior work, evidence, and motivating problem. Cite supporting sources in Appendix~E and the bibliography.

%\chapter{Stakeholder List}\ChapterLabel{StakeholderList}
%Identify users, sponsors, customers, operators, maintainers, regulators, and other affected parties. State each stakeholder's role and interest.

%\chapter{Analysis of Stakeholder Needs}\ChapterLabel{StakeholderNeeds}
%Document stakeholder needs, priorities, conflicts, and the evidence used to understand them. Link needs to requirements in \ChapterRef{Requirements}.

%\chapter{Project Scope and Resources}\ChapterLabel{ScopeResources}
%Define in-scope and out-of-scope work, team capabilities, facilities, tools, funding, and other resources required for delivery.


%Describe the work breakdown, responsibilities, schedule, risk management, review gates, and change-control approach. Include the team organization chart in Appendix~A and the Gantt chart in Appendix~B.


